\documentclass[a4paper,10.5pt]{article}
\usepackage[utf8]{inputenc}
\usepackage[T1]{fontenc}
\usepackage[margin=18mm,top=16mm,bottom=16mm]{geometry}
\usepackage{enumitem}
\usepackage{color}
\usepackage[colorlinks=true,linkcolor=black,urlcolor=blue]{hyperref}
\usepackage{array}
\usepackage{amssymb}
\usepackage{fancyhdr}
\pagestyle{fancy}
\fancyhf{}
\lhead{\small CASS 15 first-audit readiness checklist}
\rhead{\small David H. Lee \,\textperiodcentered\, ACMA CGMA \,\textperiodcentered\, September 2026}
\cfoot{\small\thepage}
\renewcommand{\headrulewidth}{0.3pt}
\setlength{\parskip}{4pt}
\setlength{\parindent}{0pt}
\setlist[itemize]{topsep=1pt,itemsep=1pt,parsep=0pt,leftmargin=*}
\setlist[enumerate]{topsep=1pt,itemsep=1pt,parsep=0pt,leftmargin=*}
\newcommand{\q}[1]{\item[$\square$] #1}

\begin{document}

{\LARGE\bfseries CASS 15 first-audit readiness checklist}\\[1mm]
{\large For UK payment and e-money firms, four months into the supplementary safeguarding regime}

\vspace{2mm}
\textbf{Why this exists.} The supplementary safeguarding regime (FCA PS25/12; CASS 15, SUP 3A, SUP 16.14A) has applied since 7 May 2026 to every authorised payment institution, authorised e-money institution, small EMI and e-money-issuing credit union. Your first safeguarding audit period is running now. The auditor reports to the FCA and opines on whether adequate systems were maintained \emph{throughout} the period. The FRC's interim guidance and the auditors' own briefings say the same thing: a control reconstructed after year end does not count. This checklist is the set of questions an auditor will effectively ask, phrased so that a CFO or designated safeguarding individual can answer them without a compliance adviser in the room. Tick what you can evidence today, not what you intend to do.

\section*{1. Scope and period}
\begin{itemize}
\q{We have confirmed which entity or entities are in scope and whether any qualifies for the audit exemption (relevant funds under \pounds100{,}000 at all times over a period of at least 53 weeks; falling below the threshold at period end is not enough).}
\q{We have agreed the first audit period with our auditor (not more than 53 weeks) and diarised the report deadline: six months after the first period ends, four months for subsequent periods.}
\q{We have appointed, or have a signed engagement with, an auditor who has confirmed capacity and CASS 15 competence for our reporting date.}
\q{One named director or senior manager is the designated safeguarding individual, and the board minutes show the appointment and the reporting line.}
\end{itemize}

\section*{2. Daily internal reconciliation}
\begin{itemize}
\q{On every reconciliation day since 7 May 2026 (business days; weekends and UK bank holidays excluded), we compared the relevant funds that should be safeguarded with our own records, at a defined cut-off time.}
\q{The cut-off time is documented and reflects our operating model, and the same cut-off is used every day.}
\q{For each day we can produce the source extracts, the calculation, the result, the preparer, the reviewer, and the time it was done. All of these were created on the day, not later.}
\q{Population completeness is checked: we can show that every client balance and every in-flight transaction was included, not just that the totals agree.}
\end{itemize}

\section*{3. Daily external reconciliation}
\begin{itemize}
\q{On every reconciliation day we compared our internal record of safeguarded funds with the statements of every bank, custodian and insurer or guarantor that holds or backs them.}
\q{Timing differences (bank cut-off versus our cut-off, FX settlement, card scheme settlement) are itemised each day, not netted into a single unexplained figure.}
\q{Where a partner's statement was late or unavailable, the exception is recorded on the day with what we did instead.}
\end{itemize}

\section*{4. Segregation position and D+1}
\begin{itemize}
\q{Each day we calculate what should be held in safeguarding accounts and compare it with what is actually held, and any shortfall is topped up from our own funds by the end of the next business day, with the transfer evidenced.}
\q{Any excess of our own money in the safeguarding accounts is identified and withdrawn within the permitted window, with the reason and the approval recorded.}
\q{Unallocated funds are allocated to the correct client by the end of the next business day; unidentified funds are recorded separately, investigated, and returned to sender if they cannot be identified.}
\end{itemize}

\section*{5. Breaks and breaches}
\begin{itemize}
\q{Every reconciliation difference is logged with an owner, the date raised, the root cause, and an ageing that is reviewed at least weekly.}
\q{Resolution of a break requires a second person's approval (maker-checker), and the approval is recorded with the break.}
\q{We have a written threshold for what constitutes a reportable breach, and every breach since May has been notified to the FCA in the monthly return with the date it occurred and the date it was fixed.}
\q{Nobody can edit a past day's reconciliation output without leaving a trace. If the answer is ``it is a spreadsheet'', the answer to this question is no.}
\end{itemize}

\section*{6. Monthly safeguarding return}
\begin{itemize}
\q{We have filed every return due since May within 15 business days of month end, and we have the submission receipts.}
\q{The figures in each return tie back to the daily reconciliations for that month, and we can show the tie-out.}
\q{Client counts, safeguarding methods, balance breakdowns, and any use of non-standard reconciliation are consistent month to month or the change is explained.}
\end{itemize}

\section*{7. Records, acknowledgement letters and third parties}
\begin{itemize}
\q{Every safeguarding account has an acknowledgement letter in the FCA-prescribed wording or a compliant pre-May letter, and each letter has been reviewed within the last year.}
\q{Our selection of each bank, custodian and insurer is documented with the due diligence performed, and a periodic review date is set.}
\q{We have assessed and documented whether we should diversify safeguarding providers, including the reasoning if we chose not to.}
\q{Any safeguarding insurance or guarantee pays directly into a safeguarding account and cannot be cancelled without the required notice to us and the FCA.}
\end{itemize}

\section*{8. Resolution pack}
\begin{itemize}
\q{We have rehearsed retrieving the resolution pack and produced it in full within 48 hours, and we recorded the time it took.}
\q{The pack identifies group members, third parties, critical staff, standard client terms and the location of every record an administrator would need.}
\q{Inaccuracies found in the pack are corrected within five business days and the correction is dated.}
\end{itemize}

\section*{9. Governance and evidence}
\begin{itemize}
\q{The designated individual has reported to the board on safeguarding at least once since May, and the paper and the minute exist.}
\q{We can hand the auditor an indexed catalogue of evidence, by requirement, by day, without building it from scratch after the request.}
\q{Someone other than the person who runs the reconciliation has tested a sample of days since May against this checklist and written down what they found.}
\end{itemize}

\vspace{2mm}
\hrule
\vspace{2mm}
\textbf{Scoring.} Count the unticked boxes in sections 2 to 6. Zero: prepare the evidence catalogue and you are ready. One to four: the gaps are fixable in weeks, but every month they remain unfixed is a month the auditor will sample. Five or more: the control is probably not operating as the rules define it, and the priority is to make it operate and evidence it from today, then address the historic period with your auditor.

\textbf{About the author.} David Lee is a CIMA-qualified finance lead with a Computer Science degree. He owned the safeguarding reconciliation at Apron Payments (FCA-authorised EMI, ClearBank), where he built the daily reconciliation in Python and cut it from a full Financial Controller day to 30 minutes before handing it over as business as usual, and he led the finance side of the accounting-rules-engine programme at Revolut that resolved a qualified audit. He does not audit and does not write safeguarding policy. He builds the control and the evidence. \href{mailto:lee.hantae@gmail.com}{lee.hantae@gmail.com} \,\textperiodcentered\, 

{\footnotesize This checklist is a practitioner's summary of PS25/12 and the FRC's interim guidance as at September 2026. It is not legal or regulatory advice. Where it conflicts with your auditor's or the FCA's stated requirements, follow them.}

\end{document}
